\documentclass[11pt,a4paper]{article}
\usepackage[margin=2.2cm]{geometry}
\usepackage[T1]{fontenc}
\usepackage{amsmath,amssymb,xcolor,enumitem,booktabs,graphicx,tabularx}
\usepackage[normalem]{ulem}
\setlength{\parskip}{0.6em}\setlength{\parindent}{0pt}
\definecolor{newc}{HTML}{1f6f43}
\newcommand{\new}[1]{\textcolor{newc}{#1}}
\newcommand{\old}[1]{\textcolor{gray}{\sout{#1}}}
\newcommand{\note}[1]{{\small\itshape [#1]}}
\newcommand{\lab}[1]{\textsf{\footnotesize\detokenize{#1}}}
\renewcommand{\ref}[1]{\lab{#1}}
\newcommand{\nref}[1]{(\lab{#1})}
\newcommand{\citep}[1]{[\lab{#1}]}
\newcommand{\citet}[1]{\lab{#1}}
\newcommand{\mockup}[1]{\fbox{\small mockup: #1}}
\title{Proposed change, the lock test}
\date{25 November 2026, version 3}
\begin{document}\sloppy
\maketitle

\section*{What changed in version 3}
\begin{itemize}[nosep]
  \item New part 3, how an episode ends. The code cuts a training episode
    and lets the critic's value stand in for the rest, and no run is ever
    stopped on a jam. The thesis says neither. Three sentences, one in the
    Method, one in the Setup and one in the threats.
\end{itemize}

\section*{What changed in version 2}
\begin{itemize}[nosep]
  \item No new runs. The text says plainly what the lock measure can and
    cannot show, instead of testing more thresholds (2a, 2h).
  \item The threshold gets its reason. Token Passing never deadlocks, so its
    longest wait is normal waiting, at most 74 steps. 200 is almost three
    times that, and it was fixed before the learned runs were counted.
  \item Question 1 is answered as option (a). The text says only what the
    data shows, that locks are rare and that decentralised agents stand
    still alone for long. The cause stays open (2e, 2f, 2g).
  \item Appendix L gives the counts at 200 and at 500 steps, both from the
    replays (2i).
  \item The served table at 16 agents is left to its own review (question 4).
\end{itemize}

\section*{How to read this}
The lock is now measured as agreed on 25 November. Two agents on
neighbouring vertices both stand still, each holding its task, for 200 steps
in a row. All 70 episodes that could reach 200 steps were replayed, and every
replay gave its original row back. This review puts the new counts in the
thesis and changes every sentence that leans on the old count.

Black text stays, \new{green} is new and \old{struck} text goes.
\note{Notes in brackets} say where a number comes from. Nothing below is in
the thesis yet. The questions are at the end.

\textbf{Overlap.} \lab{review_locks_fleet_reward} (version 4) already holds
the new lock definition (its part 4.0) and leaves the lock counts as
mockups. This review fills them. Once both are agreed they are applied
together, the lock test as [A136] and the fleet reward test as [A143].

\textbf{The main finding changes the argument.} Locks are much rarer than
the old count said. Where they do happen, the seeds that lock serve no less
than the seeds that do not. So locks do not explain why the learned
controllers fall behind, and the text says they do in five places (parts 2c,
2e, 2f, 2g and 2h). Question 1 asks what to say instead.

\section*{Part 1. What the replays show}

\textbf{16 agents.} Seeds that lock, out of 20, each rule played by the
policy trained under it. In brackets the old count (two agents hold a task
for 500 steps), as \lab{tab:exp:locks} has it now.

\begin{center}\small
\begin{tabular}{l ccc}
  \toprule
  controller & fixed & demand & congestion-aware \\
  \midrule
  Token Passing & 0 (0) & 0 (0) & 0 (0) \\
  centralised   & 0 (5) & \mockup{eval-pc} (3) & 0 (0) \\
  section-based & 0 (5) & 0 (8) & 0 (3) \\
  decentralised & 9 (20) & 1 (13) & 0 (3) \\
  with messages & 2 (15) & 0 (20) & 1 (18) \\
  \bottomrule
\end{tabular}
\end{center}

\note{Token Passing's longest stand-still of a single agent is 30 steps at 16
agents, so no pair can reach 200. Centralised demand is still being tested
on the PC. Its test rows are empty, so its count follows.}

\textbf{Locked seeds serve no less.} Decentralised under fixed storage locks
on 9 seeds. They serve 0.19 to 0.56, mean 0.34. The 11 seeds that do not lock
serve 0.06 to 0.41, mean 0.24. With messages, the two locked seeds under fixed
storage serve 0.51 and 0.68 against a mean of 0.63. Only one locked seed is
far below the rest (with messages, congestion-aware, 0.08).

\textbf{What the shortfall goes with instead.} The longest time a single
agent stands still, median over the 20 seeds, is 9 to 16 steps for the
centralised controller and 16 to 78 for the section-based one. For the
decentralised controller it is 80 to 321 steps, and with messages 47 to 118.
So decentralised agents stand still alone for a long time, without a
neighbour that stands still too. \note{From the test rows,
\lab{longest_stand_still}. Not yet checked where these agents stand.}

\textbf{8, 32 and 64 agents} (the 16-agent congestion-aware policies, test
seeds). Seeds that lock, out of 20.

\begin{center}\small
\begin{tabular}{l cccc}
  \toprule
  controller & 8 & 16 & 32 & 64 \\
  \midrule
  centralised   & 0 & 0 & 0 & 0 \\
  section-based & 0 & 0 & 1 & 3 \\
  decentralised & 0 & 0 & 3 & 7 \\
  with messages & 0 & 1 & 7 & 7 \\
  Token Passing & 0 & 0 & 0 & 0 \\
  \bottomrule
\end{tabular}
\end{center}

\note{Token Passing at 32 and 64 from its test-seed rows, which carry the new
measure. Its longest pair stand-still is 31 and 66 steps, its longest single
stand-still 43 and 74.}

\textbf{The layout.} Under congestion-aware storage, up to 32 agents every lock of the section-based and
decentralised controllers sits on a seed where one category has all its
workstations on one-lane aisles. With messages also locks on 1 other seed at
32 agents, and at 64 agents decentralised locks on 4 other seeds and with
messages on 2. At 16 agents decentralised under fixed storage locks mostly
off these seeds (8 of its 9 locks). \note{\lab{locks_README.md}, steps 2 and
2b. Where the locks sit, from the position trace, follows in a mockup
below.}

\section*{Part 2. Proposed text}

\subsection*{2a. Setup, where a lock is defined (5.Experiments, before
\lab{tab:exp:locks})}
\note{The definition is part 4.0 of the fleet reward review. Repeated here
with the numbers filled.}

\old{Locks explain most of the decentralised controller's failure at 16
agents (Table~\ref{tab:exp:locks}). A seed locks when two agents hold the
same task for 500 steps or more. The decentralised controller locks on 13 to
20 seeds under fixed and demand storage, and with messages on 15 under fixed
storage. Token Passing never locks.}

\new{A seed locks when two agents on neighbouring vertices both stand still,
each holding its task, for 200 steps in a row. The threshold is set high on
purpose. The measure only watches and never stops a run, so two agents that
wait and then clear the way themselves do not count as a lock. It is set from
Token Passing, which never deadlocks in these warehouses. Its agents wait at
most 74 steps, at any fleet size, so that is how long normal waiting lasts. The
threshold is almost three times as long and was fixed before the learned runs
were counted.}

\new{A lock is a symptom with several causes. The layout, a controller's choice
to wait and the queue behind the two agents all add to it, and the count
cannot tell them apart. Stops shorter than 200 steps are not counted as locks
and show up in the blocked time instead.}

\new{Locks are rare at 16 agents (Table~\ref{tab:exp:locks}). The centralised
and section-based controllers never lock. The decentralised controller locks
on 9 seeds under fixed storage and on 1 under demand storage. With messages it
locks on at most 2 seeds. Token Passing never locks.}

\new{Locks do not explain why a controller falls behind. Under fixed storage
the 9 seeds where the decentralised controller locks serve 0.34 on average,
and the 11 where it does not serve 0.24.}

\note{\lab{tab:exp:locks} takes the counts of part 1. The two TODO comments
above it go. Caption unchanged.}

\subsection*{2b. ``Why the count does not follow the mean'' and
\lab{fig:exp:seeds}}

\old{Locks spread the seeds further. When two agents hold one task for
hundreds of steps, the work behind them waits, and the seed falls far below
the rest. The section-based controller under fixed storage locks on five
seeds, and four of them serve only 0.39 to 0.53.}

\note{The paragraph goes. Section-based under fixed storage locks on no seed
under the new measure.}

Caption: Served share of each test seed at 16 agents. \old{A seed that locks
can fall far below the rest.}

\note{\lab{fleet16_seeds.py} marks a locked seed with a cross by the old
count. Redraw with the replay counts, or drop the crosses. With 13 locked
seeds in all, the crosses would show that locked seeds are not the low
ones.}

\subsection*{2c. ``RQ4, communication'', the end}
\old{The decentralised controller with messages still locks on 15 to 20 seeds,
so messages do not let two agents back out of a lock.}
\new{Messages make locks rarer. Under fixed storage the decentralised
controller locks on 9 seeds, and with messages on 2.}

\note{Fixed storage is the only rule where the plain decentralised controller
locks often, so it is the only fair comparison. Part 4b of the fleet reward
review (``the same holds at 32 and 64'') no longer holds. At 32 agents with
messages locks more (7 against 3).}

\subsection*{2d. ``RQ3, scale'', after ``none keeps up on any seed''}
\new{Locks appear as the fleet grows. At 8 and 16 agents the
congestion-aware policies do not lock. At 64 agents the decentralised
controller locks on 7 seeds, with messages on 7 and the section-based
controller on 3. The centralised controller never locks. Token Passing locks
on no seed at 32 or 64 agents.}

\new{\mockup{One sentence on where the locks sit, from the position trace
(step 3, running on the Lenovo).}}

\note{This fills the lock mockup of part 4a of the fleet reward review. Its
table \lab{tab:exp:fleets} gets the lock column of part 1.}

\subsection*{2e. Discussion, ``Why the learned controllers fall behind''}
\old{The learned controllers fall behind mostly through locks. Two agents that
each wait for the other hold one task for the rest of the run, and the agents
behind them queue. \citet{krnjaic2022warehouse_marl} saw the same with
agents that share one policy, and a manager that spread the goals over the
warehouse avoided it. The section-based controller does much the same and
locks less.}

\new{The learned controllers do not fall behind through locks. Locks are rare,
and the seeds that lock serve no less than the others. What sets the
decentralised controllers apart is that a single agent stands still for a
long time while it holds a task, with no stuck agent beside it. Its longest
stand-still is 80 to 321 steps on a typical seed, against 9 to 16 for the
centralised controller. These waits add to the blocked time, which is 13 to
26 steps per pick for the learned controllers and under 2 for Token Passing.
Why the agents wait this long was not traced.}

\note{The second sentence of the old text (\citet{krnjaic2022warehouse_marl})
goes with it, since it explained locks. The rest of the paragraph (episode length, the long runs) stays.}

\subsection*{2f. ``What could have been done differently'', the item on locks}
\old{\emph{A way out of locks.} Most of the decentralised controller's loss is
locks. A shared credit, or a rule that makes one of two locked agents give
way, would let them back out.}

\new{\emph{Where agents wait.} The decentralised agents stand still alone for
long stretches. A trace of where they stand and what their goal is would show
whether they wait before goals that cost them, or for another reason.}

\note{The fleet reward test still asks whether shared credit helps, but no
longer as a cure for locks.}

\subsection*{2g. Related work, the end of the van Knippenberg sentence}
Here waiting is the cheaper choice\old{, which helps explain the locks}.
\new{, which may be why decentralised agents stand still for so long}.

\subsection*{2h. Threats to validity, the item on one-lane aisles}
Agents can lock at a workstation on an aisle one lane wide
(Section~\ref{sec:method:resolution}). A category with all its workstations
on such aisles occurs in 14\,\% of the warehouses at 8 agents and 27\,\% at
32 agents, and in 5 of the 20 test warehouses at 8 agents. \old{At 16 agents
the decentralised controller locks on every test seed under fixed and demand
storage. Token Passing does not lock, so in these warehouses the gap between
it and the learned controllers is partly the layout.} \new{Locks are rare,
and under congestion-aware storage up to 32 agents those of the
section-based and decentralised controllers sit on such warehouses. Locked seeds serve no less than the rest, so the
layout explains little of the gap to Token Passing. The lock measure itself
is limited. It sees two neighbouring agents that stand still for 200 steps,
not why they stand still, and a jam of three agents or a stop of 150 steps
does not count.} A shared reward for the
decentralised controller, which would give it a reason to back out, is left
to future work.

\note{The last sentence changes once the fleet reward test is in.}

\subsection*{2i. Appendix L, the lock paragraph and \lab{tab:app:locks}}
\old{A seed locks when two agents each hold their own task for 500 steps or
more at the same time. Table~\ref{tab:app:locks} counts these seeds at both
fleets, with the stricter threshold of 1250 steps beside it. Fewer seeds lock
at 8 agents, at most 10 of 20. At 16 agents the decentralised controller locks
on every seed under fixed storage, and on all 20 even at 1250 steps.}

\new{Table~\ref{tab:app:locks} counts the seeds that lock, as defined in
Section~\ref{sec:exp:results}, with the longest lock of each run beside it.
At 8 agents no policy comes near 200 steps. At 16 agents only the two
decentralised controllers lock. Under fixed storage the decentralised
controller locks on 9 seeds at 200 steps and on 2 at 500, and its longest lock
lasts 1859 steps.}

\note{Table: one row per controller, 8 agents congestion-aware only (all 0), 16
agents each rule, cells ``200 / 500''. From the replays. 16 agents: decentralised
fixed 9 / 2, demand 1 / 1, congestion-aware 0 / 0. With messages fixed 2 / 0,
demand 0 / 0, congestion-aware 1 / 0. Centralised and section-based 0 / 0,
centralised demand to follow.}

\note{Caption: ``Test seeds that lock, out of 20, at 200 and at 500 steps.
Token Passing locks on no seed.''}

\subsection*{2j. Method, \lab{sec:method:resolution}}
\note{No change proposed. It describes how two agents can lock at a
one-lane workstation and why no rule breaks it, which the counts confirm
for the section-based controller. It predicts no rates. Its sentence ``A
decentralised agent \ldots gains nothing when it backs out'' fits the
counts, since the decentralised controller is the one that locks.}

\section*{Part 3. How an episode ends}

\note{Checked in the code. RLlib marks the end of a training episode as
truncated, not terminated (\lab{multi_agent_env.py}, \lab{training/env.py}),
so PPO bootstraps the value of the last state. No run is stopped when agents
jam, the reason is in GAPS.md, the stall rule that was not adopted.}

\subsection*{3a. Method, \lab{sec:method:training}, a new paragraph after the
first (``\ldots are in Appendix~\ref{app:training}.'')}

\new{A training episode is cut at its length, not ended. The value estimate of
the last state stands in for everything after the cut. So the policy is not
taught that the warehouse stops there, and work it leaves open still costs
what the value estimate expects it to cost later.}

\subsection*{3b. Setup, ``Training'', after ``\ldots so a policy cannot learn
one map by heart.''}

\new{No run is stopped when agents jam. In an early run a policy recovered
from stops of 94 and 223 steps, so a stop would have cut runs that recover. A
stop would also make a jam cheaper, because the penalty for a refused move ends
with the run. A jam is therefore seen in full, in training and in the test.}

\note{The 94 and 223 steps are in GAPS.md with the stall rule. This is a reason
for a design choice, so it stays in the Setup and not in the Results.}

\subsection*{3c. Threats to validity, a new item after the one on checkpoints}

\begin{itemize}\item \new{No cost is put on work left open when a training episode is cut.
The policy only sees it through its value estimate, so a policy whose estimate
is poor may care too little about work it leaves late in an episode. Training
episodes of 16\,000 steps, longer than a test run, make this smaller. In the
test, work left open counts as not served.}
\end{itemize}

\section*{Questions}
\begin{enumerate}
  \item \textbf{What replaces ``fall behind mostly through locks''?} Two
    options.
    (a) Say only what the data shows. Locks are rare and do not explain
    the shortfall, and the decentralised agents stand still alone for long
    stretches. The cause is left open.
    (b) First check where those lone agents stand and why, with the same
    position trace as step 3, for a few decentralised seeds at 16 agents
    (about 1 hour on the Lenovo). The RL review found that agents stop
    before goals that bring a penalty, which may be the same thing.
    I recommend (b), since the Discussion needs a reason and (a) leaves a
    gap where the old one was.
  \item \textbf{8 agents in \lab{tab:app:locks}.} The 8-agent column comes
    from the runs trained at 8 agents. Seven of their episodes can reach 200
    steps (decentralised fixed 3 and congestion-aware 2, with messages
    demand 2), and they would need a replay on the code that evaluated
    them. Or the column takes the scale test at 8 agents, where nothing
    locks, as Appendix L already does for its other 8-agent tables
    ([A142]). I recommend the scale test.
  \item \textbf{The old measure.} Drop it everywhere, including the 1250
    column and the cross in \lab{fig:exp:seeds}? I recommend yes. It counts
    any two slow tasks anywhere and grows with the fleet without a jam.
  \item \textbf{The served table.} \lab{tab:exp:served} still has first-run
    numbers under its TODO (for example centralised fixed 0.54, now 0.78
    from the long run). Not part of this review, but 2a quotes served
    shares, so the two should be updated together.
\end{enumerate}

\section*{Answers so far (25 November)}
\begin{itemize}[nosep]
  \item 2. The scale test at 8 agents, checked true. Appendix L plays only
    the congestion-aware policies there, and none of their 20 seeds comes
    near 200 steps. The 8-agent part of \lab{tab:app:locks} keeps the
    congestion-aware column only, all 0.
  \item 3. The old measure is dropped everywhere.
  \item 1. Option (a), no extra trace. The text states the limits of the
    measure instead of testing more thresholds (author: no time for more
    runs).
  \item 4. Left to its own review. The long runs change more than numbers.
    With messages serves 0.69 under demand storage (was 0.28) and 0.78 under
    congestion-aware storage (was 0.34), so the RQ4 sentence ``under demand
    storage messages lower the share'' turns around. That review follows
    when centralised demand is tested.
\end{itemize}

\section*{Still to do before applying}
\begin{itemize}[nosep]
  \item Centralised demand at 16 agents, once eval-pc writes its test rows.
  \item Where the locks sit (step 3, Lenovo, after Token Passing).
  \item Apply with the fleet reward review, [A136] in GAPS.md (lock test and
    part 3).
\end{itemize}

\end{document}
